408
\hfill 6 Integration
b) if the limit in a) is equal to $1$ (and it is), then
$$
\text{erf}(x) = \sqrt{\frac{2}{\pi}} \int_{0}^{x} e^{-t^2} dt = \sqrt{\frac{2}{\pi}} \left( \frac{\sqrt{\pi}}{2} \left( 1 - \frac{2}{\sqrt{\pi}} \int_{x}^{+\infty} e^{-t^2} dt \right) \right)
$$
$$
= 1 - \frac{2}{\sqrt{\pi}} e^{-x^2} \left( \frac{1}{2x} - \frac{1}{2^2 x^3} + \frac{1 \cdot 3}{2^3 x^5} - \frac{1 \cdot 3 \cdot 5}{2^4 x^7} + o\left(\frac{1}{x^7}\right) \right)
$$
as $x \to +\infty$.

\textbf{6.} Prove the following statements.
a) If a heavy particle slides under gravitational attraction along a curve given
in parametric form as $x = x(\theta), y = y(\theta)$, and at time $t = 0$ the particle had zero
velocity and was located at the point $x_0 = x(\theta_0), y_0 = y(\theta_0)$, then the following
relation holds between the parameter $\theta$ defining a point on the curve and the time $t$
at which the particle passes this point (see formula (6.63) of Sect. 6.4)
$$
t = \pm \int_{\theta_0}^{\theta} \sqrt{\frac{(x'(\theta))^2 + (y'(\theta))^2}{2g(y_0 - y(\theta))}} d\theta,
$$
in which the improper integral necessarily converges if $y'(\theta_0) \neq 0$. (The ambiguous
sign is chosen positive or negative according as $t$ and $\theta$ have the same kind of
monotonicity or the opposite kind; that is, if an increasing $\theta$ corresponds to an
increasing $t$, then one must obviously choose the positive sign.)

b) The period of oscillation of a particle in a well having cross section in the
shape of a cycloid
$$
x = R(\theta + \pi + \sin\theta), \quad |\theta| \le \pi,
$$
$$
y = -R(1 + \cos \theta),
$$
is independent of the level $y_0 = -R(1 + \cos \theta_0)$ from which it begins to slide and is
equal to $4\pi \sqrt{R/g}$ (see Problem 4 of Sect. 6.4).